\begingroup\setbox0=\vbox{\section*{Hình tam giác. Diện tích hình tam giác}}\endgroup
\chapterhead{11}{Hình tam giác. Diện tích hình tam giác}

\shead{Đáy, chiều cao và diện tích}
\begin{knowbox}
\begin{minipage}[t]{0.6\linewidth}
\textbf{Hình tam giác} có $3$ cạnh, $3$ đỉnh và $3$ góc. Theo góc, có hình tam giác \emph{nhọn} (ba góc nhọn), hình tam giác \emph{tù} (một góc tù) và hình tam giác \emph{vuông} (một góc vuông).\par
\smallskip
Đoạn thẳng kẻ từ một đỉnh và \textbf{vuông góc} với cạnh đối diện là \textbf{đường cao}; cạnh đó là \textbf{cạnh đáy}. Mỗi hình tam giác có $3$ đường cao.\par
\smallskip
\textbf{Diện tích hình tam giác} bằng độ dài đáy nhân với chiều cao (cùng một đơn vị đo) rồi chia cho $2$:
\[S=\frac{a\times h}{2}\qquad h=\frac{S\times2}{a}\qquad a=\frac{S\times2}{h}\]
Trong hình tam giác vuông, hai cạnh góc vuông là đáy và chiều cao của nhau.
\end{minipage}\hfill
\begin{minipage}[t]{0.36\linewidth}\vspace{2pt}
\centering
\begin{tikzpicture}[scale=0.9]
\coordinate (B) at (0,0); \coordinate (C) at (4,0); \coordinate (A) at (1.3,2.6); \coordinate (H) at (1.3,0);
\fill[tmBlue!12] (A)--(B)--(C)--cycle;
\draw[geo] (A)--(B)--(C)--cycle;
\draw[dashed,line width=0.7pt,tmOrange!80!black] (A)--(H);
\ramark{H}{C}{A}
\node[lbl,above] at (A){A}; \node[lbl,below left] at (B){B}; \node[lbl,below right] at (C){C}; \node[lbl,below] at (H){H};
\node[dim,right] at (1.3,1.3){$h$};
\draw[|<->|,Slate] (0,-0.55)--node[fill=white,dim]{$a$}(4,-0.55);
\end{tikzpicture}
\end{minipage}
\end{knowbox}
\gap
\begin{exbox}{\ExWord\ 1}
Tính diện tích hình tam giác có:
\begin{cols2}
\pa{a} độ dài đáy $8$ cm, chiều cao $5$ cm; & \pa{b} độ dài đáy $2{,}5$ dm, chiều cao $12$ cm.
\end{cols2}
\guide
\pa{a} $S=8\times5:2=20$ cm$^2$.\par
\pa{b} Trước hết đổi về cùng đơn vị đo: $2{,}5$ dm $=25$ cm.\par $S=25\times12:2=150$ cm$^2$.
\dapso{\pa{a} $20$ cm$^2$;\par \pa{b} $150$ cm$^2$.}
\end{exbox}
\gap
\begin{exbox}{\ExWord\ 2}
\begin{minipage}[t]{0.62\linewidth}
Hình tam giác ABC vuông tại A có $AB=6$ cm, $AC=8$ cm và $BC=10$ cm. Tính diện tích hình tam giác đó và độ dài đường cao AH.
\end{minipage}\hfill
\begin{minipage}[t]{0.34\linewidth}\vspace{0pt}\centering
\begin{tikzpicture}[scale=0.36]
\coordinate (A) at (0,0); \coordinate (B) at (0,6); \coordinate (C) at (8,0); \coordinate (H) at (2.88,3.84);
\draw[geo] (A)--(B)--(C)--cycle; \ramark{A}{B}{C}
\draw[dashed,line width=0.7pt,tmOrange!80!black] (A)--(H); \ramark{H}{C}{A}
\node[lbl,below left] at (A){A}; \node[lbl,above] at (B){B}; \node[lbl,below right] at (C){C}; \node[lbl,above right] at (H){H};
\end{tikzpicture}
\end{minipage}
\guide
Lấy AC làm đáy thì AB là chiều cao: $S=6\times8:2=24$ cm$^2$.\par
Lấy BC làm đáy thì AH là chiều cao: $AH=24\times2:10=4{,}8$ cm.
\dapso{$24$ cm$^2$;\par $AH=4{,}8$ cm.}
\end{exbox}
\gap
\begin{exbox}{\ExWord\ 3}
Một hình tam giác có diện tích $36$ cm$^2$ và độ dài đáy $9$ cm. Tính chiều cao của hình tam giác đó.
\guide
$h=S\times2:a=36\times2:9=8$ cm.
\dapso{$8$ cm.}
\end{exbox}
\gap
\begin{exbox}{\ExWord\ 4}
\begin{minipage}[t]{0.6\linewidth}
Hình chữ nhật ABCD có $AB=12$ cm và $BC=8$ cm; M là trung điểm của BC. Tính diện tích các hình tam giác ABM, MCD và AMD.
\end{minipage}\hfill
\begin{minipage}[t]{0.36\linewidth}\vspace{0pt}\centering
\begin{tikzpicture}[scale=0.3]
\coordinate (A) at (0,8); \coordinate (B) at (12,8); \coordinate (C) at (12,0); \coordinate (D) at (0,0); \coordinate (M) at (12,4);
\fill[shade] (A)--(M)--(D)--cycle;
\draw[geo] (A)--(B)--(C)--(D)--cycle; \draw[geo] (A)--(M)--(D);
\node[lbl,above left] at (A){A}; \node[lbl,above right] at (B){B}; \node[lbl,below right] at (C){C}; \node[lbl,below left] at (D){D}; \node[lbl,right] at (M){M};
\end{tikzpicture}
\end{minipage}
\guide
$BM=MC=8:2=4$ cm.\par $S_{ABM}=12\times4:2=24$ cm$^2$;\par $S_{MCD}=12\times4:2=24$ cm$^2$.\par
$S_{ABCD}=12\times8=96$ cm$^2$, nên $S_{AMD}=96-24-24=48$ cm$^2$.\par
\emph{Nhận xét:} hình tam giác AMD có đáy AD và chiều cao bằng AB, nên diện tích của nó bằng một nửa diện tích hình chữ nhật.
\dapso{$24$ cm$^2$; $24$ cm$^2$; $48$ cm$^2$.}
\end{exbox}
\dividerbar{Luyện tập}
\begin{probbox}{1}{\relax}
Tính diện tích hình tam giác có độ dài đáy $a$ và chiều cao $h$:
\begin{cols2}
\pa{a} $a=15$ cm, $h=8$ cm; & \pa{b} $a=4{,}5$ m, $h=2$ m;\\
\pa{c} $a=\dfrac34$ m, $h=\dfrac25$ m; & \pa{d} $a=1{,}2$ dm, $h=15$ cm.
\end{cols2}
\end{probbox}
\probsol{
\pa{a} $15\times8:2=60$ cm$^2$.\par \pa{b} $4{,}5\times2:2=4{,}5$ m$^2$.\par \pa{c} $\dfrac34\times\dfrac25:2=\dfrac3{20}$ m$^2$.\par \pa{d} $1{,}2$ dm $=12$ cm; $12\times15:2=90$ cm$^2$.
\dapso{$60$ cm$^2$; $4{,}5$ m$^2$; $\dfrac3{20}$ m$^2$; $90$ cm$^2$.}}
\gap
\begin{probbox}{2}{\relax}
Tính diện tích hình tam giác vuông có hai cạnh góc vuông dài $3{,}5$ dm và $2$ dm.
\end{probbox}
\probsol{
$S=3{,}5\times2:2=3{,}5$ dm$^2$.
\dapso{$3{,}5$ dm$^2$.}}
\gap
\begin{probbox}{3}{\relax}
Một hình tam giác có diện tích $45$ cm$^2$ và chiều cao $6$ cm. Tính độ dài đáy của hình tam giác đó.
\end{probbox}
\probsol{
$a=45\times2:6=15$ cm.
\dapso{$15$ cm.}}
\gap
\begin{probbox}{4}{\relax}
Một mảnh đất hình tam giác có độ dài đáy $40$ m, chiều cao bằng $\dfrac34$ độ dài đáy. Tính diện tích mảnh đất đó.
\end{probbox}
\probsol{
Chiều cao: $40\times\dfrac34=30$ m.\par Diện tích: $40\times30:2=600$ m$^2$.
\dapso{$600$ m$^2$.}}
\gap
\begin{probbox}{5}{\relax}
Một tấm bìa hình tam giác có độ dài đáy $2{,}4$ dm, chiều cao bằng $\dfrac23$ độ dài đáy. Tính diện tích tấm bìa đó.
\end{probbox}
\probsol{
Chiều cao: $2{,}4\times\dfrac23=1{,}6$ dm.\par Diện tích: $2{,}4\times1{,}6:2=1{,}92$ dm$^2$.
\dapso{$1{,}92$ dm$^2$.}}
\gap
\begin{probbox}{6}{\relax}
Hình tam giác ABC vuông tại A có $AB=9$ cm, $AC=12$ cm và $BC=15$ cm. Tính độ dài đường cao AH.
\end{probbox}
\probsol{
$S=9\times12:2=54$ cm$^2$.\par $AH=54\times2:15=7{,}2$ cm.
\dapso{$7{,}2$ cm.}}
\gap
\begin{probbox}{7}{\relax}
Một hình tam giác có diện tích bằng diện tích hình vuông cạnh $6$ cm. Độ dài đáy của hình tam giác là $9$ cm. Tính chiều cao của hình tam giác đó.
\end{probbox}
\probsol{
Diện tích: $6\times6=36$ cm$^2$.\par Chiều cao: $36\times2:9=8$ cm.
\dapso{$8$ cm.}}
\gap
\begin{probbox}{8}{\relax}
Một hình tam giác có độ dài đáy $18$ cm. Nếu kéo dài đáy thêm $4$ cm (giữ nguyên chiều cao) thì diện tích tăng thêm $20$ cm$^2$. Tính diện tích hình tam giác ban đầu.
\end{probbox}
\probsol{
Phần tăng thêm là hình tam giác có đáy $4$ cm và diện tích $20$ cm$^2$; chiều cao của nó là $20\times2:4=10$ cm.\par
Diện tích ban đầu: $18\times10:2=90$ cm$^2$.
\dapso{$90$ cm$^2$.}}
\gap
\begin{probbox}{9}{\relax}
Hình vuông ABCD có cạnh $10$ cm; M là trung điểm của BC. Tính diện tích hình tam giác ABM và hình tứ giác AMCD.
\end{probbox}
\probsol{
$S_{ABM}=10\times5:2=25$ cm$^2$.\par $S_{AMCD}=10\times10-25=75$ cm$^2$.
\dapso{$25$ cm$^2$; $75$ cm$^2$.}}
\gap
\begin{probbox}{10}{\relax}
Hình tam giác ABC có diện tích $60$ cm$^2$ và $BC=15$ cm. Điểm M nằm trên BC sao cho $BM=5$ cm. Tính diện tích hình tam giác ABM.
\end{probbox}
\probsol{
Chiều cao hạ từ A xuống BC là $60\times2:15=8$ cm.\par Hình tam giác ABM có cùng chiều cao đó: $S_{ABM}=5\times8:2=20$ cm$^2$.
\dapso{$20$ cm$^2$.}}
\gap
\begin{probbox}{11}{\relax}
\begin{minipage}[t]{0.62\linewidth}
Trong hình bên có bao nhiêu hình tam giác? Hãy kể tên các hình tam giác đó.
\end{minipage}\hfill
\begin{minipage}[t]{0.34\linewidth}\vspace{0pt}\centering
\begin{tikzpicture}[scale=0.6]
\coordinate (A) at (2,2.6); \coordinate (B) at (0,0); \coordinate (M) at (1.3,0); \coordinate (N) at (2.6,0); \coordinate (C) at (4,0);
\draw[geo] (A)--(B)--(C)--cycle; \draw[geo] (A)--(M) (A)--(N);
\node[lbl,above] at (A){A}; \node[lbl,below] at (B){B}; \node[lbl,below] at (M){M}; \node[lbl,below] at (N){N}; \node[lbl,below] at (C){C};
\end{tikzpicture}
\end{minipage}
\end{probbox}
\probsol{
Mỗi hình tam giác có đỉnh A và đáy là đoạn thẳng nối hai trong bốn điểm B, M, N, C.\par Có $3+2+1=6$ đoạn thẳng như vậy.\par
Các hình tam giác: ABM, AMN, ANC, ABN, AMC, ABC.
\dapso{$6$ hình tam giác.}}
\gap
\begin{probbox}{12}{\relax}
Tính diện tích hình tam giác có độ dài đáy $14$ cm và chiều cao $9$ cm.
\end{probbox}
\probsol{
$S=14\times9:2=63$ cm$^2$.
\dapso{$63$ cm$^2$.}}
\gap
\begin{probbox}{13}{\relax}
Tính diện tích hình tam giác có độ dài đáy $2{,}5$ m và chiều cao $1{,}6$ m.
\end{probbox}
\probsol{
$S=2{,}5\times1{,}6:2=2$ m$^2$.
\dapso{$2$ m$^2$.}}
\gap
\begin{probbox}{14}{\relax}
Tính diện tích hình tam giác có độ dài đáy $\dfrac34$ m và chiều cao $\dfrac25$ m.
\end{probbox}
\probsol{
$S=\dfrac34\times\dfrac25:2=\dfrac3{20}$ m$^2$.
\dapso{$\dfrac3{20}$ m$^2$.}}
\gap
\begin{probbox}{15}{\relax}
Tính diện tích hình tam giác có độ dài đáy $3$ dm và chiều cao $25$ cm.
\end{probbox}
\probsol{
$3$ dm $=30$ cm.\par $S=30\times25:2=375$ cm$^2$.
\dapso{$375$ cm$^2$.}}
\gap
\begin{probbox}{16}{\relax}
Tính diện tích hình tam giác vuông có hai cạnh góc vuông dài $7$ cm và $10$ cm.
\end{probbox}
\probsol{
Một cạnh góc vuông là đáy, cạnh góc vuông kia là chiều cao.\par $S=7\times10:2=35$ cm$^2$.
\dapso{$35$ cm$^2$.}}
\gap
\begin{probbox}{17}{\relax}
Một hình tam giác có diện tích $54$ cm$^2$ và độ dài đáy $12$ cm. Tính chiều cao của hình tam giác đó.
\end{probbox}
\probsol{
Chiều cao $=$ diện tích $\times2\,:$ độ dài đáy.\par $h=54\times2:12=9$ cm.
\dapso{$9$ cm.}}
\gap
\begin{probbox}{18}{\relax}
Một hình tam giác có diện tích $7{,}5$ dm$^2$ và chiều cao $2{,}5$ dm. Tính độ dài đáy của hình tam giác đó.
\end{probbox}
\probsol{
Độ dài đáy $=$ diện tích $\times2\,:$ chiều cao.\par $a=7{,}5\times2:2{,}5=6$ dm.
\dapso{$6$ dm.}}
\gap
\begin{probbox}{19}{\relax}
Một thửa ruộng hình tam giác có độ dài đáy $50$ m và chiều cao $24$ m. Cứ mỗi mét vuông thu hoạch được $0{,}8$ kg rau. Hỏi trên cả thửa ruộng đó thu hoạch được bao nhiêu ki-lô-gam rau?
\end{probbox}
\probsol{
Diện tích: $50\times24:2=600$ m$^2$.\par Số rau thu hoạch được: $0{,}8\times600=480$ kg.
\dapso{$480$ kg.}}
\gap
\begin{probbox}{20}{\relax}
Một hình tam giác có độ dài đáy $16$ cm, chiều cao bằng $\dfrac34$ độ dài đáy. Tính diện tích hình tam giác đó.
\end{probbox}
\probsol{
Chiều cao: $16\times\dfrac34=12$ cm.\par Diện tích: $16\times12:2=96$ cm$^2$.
\dapso{$96$ cm$^2$.}}
\gap
\begin{probbox}{21}{\relax}
Tổng độ dài đáy và chiều cao của một hình tam giác là $30$ cm, độ dài đáy hơn chiều cao $6$ cm. Tính diện tích hình tam giác đó.
\end{probbox}
\probsol{
Độ dài đáy: $(30+6):2=18$ cm.\par Chiều cao: $12$ cm.\par Diện tích: $18\times12:2=108$ cm$^2$.
\dapso{$108$ cm$^2$.}}
\gap
\begin{probbox}{22}{\relax}
Hình tam giác ABC vuông tại A có $AB=15$ cm, $AC=20$ cm, $BC=25$ cm. Tính diện tích hình tam giác đó và độ dài đường cao AH.
\end{probbox}
\probsol{
Diện tích: $15\times20:2=150$ cm$^2$.\par Lấy BC làm đáy thì AH là chiều cao.\par $AH=150\times2:25=12$ cm.
\dapso{$150$ cm$^2$; $AH=12$ cm.}}
\gap
\begin{probbox}{23}{\relax}
Hình tam giác thứ nhất có độ dài đáy $12$ cm và chiều cao $8$ cm. Hình tam giác thứ hai có cùng diện tích và có độ dài đáy $16$ cm. Tính chiều cao của hình tam giác thứ hai.
\end{probbox}
\probsol{
Diện tích: $12\times8:2=48$ cm$^2$.\par Chiều cao của hình tam giác thứ hai: $48\times2:16=6$ cm.
\dapso{$6$ cm.}}
\gap
\begin{probbox}{24}{\relax}
Mỗi lá cờ hình tam giác có độ dài đáy $30$ cm và chiều cao $20$ cm. Hỏi làm $20$ lá cờ như thế cần bao nhiêu đề-xi-mét vuông vải? (Không tính phần đường may.)
\end{probbox}
\probsol{
Một lá cờ: $30\times20:2=300$ cm$^2$.\par $20$ lá cờ: $6\,000$ cm$^2$ $=60$ dm$^2$.
\dapso{$60$ dm$^2$.}}
\gap
\begin{probbox}{25}{\relax}
Một hình tam giác vuông có diện tích $24$ cm$^2$ và một cạnh góc vuông dài $6$ cm. Tính độ dài cạnh góc vuông còn lại.
\end{probbox}
\probsol{
Cạnh góc vuông còn lại là chiều cao ứng với đáy $6$ cm.\par Độ dài: $24\times2:6=8$ cm.
\dapso{$8$ cm.}}
\gap
\begin{probbox}{26}{\relax}
Hình chữ nhật ABCD có chiều dài $12$ cm, chiều rộng $8$ cm. Đường chéo AC chia hình chữ nhật thành hai hình tam giác. Tính diện tích mỗi hình tam giác.
\end{probbox}
\probsol{
Hình tam giác ABC vuông tại B có hai cạnh góc vuông dài $12$ cm và $8$ cm.\par $S_{ABC}=12\times8:2=48$ cm$^2$.\par Tương tự, $S_{ACD}=48$ cm$^2$ (bằng một nửa diện tích hình chữ nhật).
\dapso{Mỗi hình $48$ cm$^2$.}}
\gap
\begin{probbox}{27}{\relax}
Một hình tam giác có diện tích $150$ cm$^2$. Nếu tăng chiều cao thêm $3$ cm (giữ nguyên đáy) thì diện tích tăng thêm $30$ cm$^2$. Tính độ dài đáy và chiều cao của hình tam giác đó.
\end{probbox}
\probsol{
Phần tăng thêm là hình tam giác có cùng đáy, chiều cao $3$ cm và diện tích $30$ cm$^2$.\par Độ dài đáy: $30\times2:3=20$ cm.\par Chiều cao: $150\times2:20=15$ cm.
\dapso{Đáy $20$ cm, chiều cao $15$ cm.}}
\gap
\begin{probbox}{28}{\relax}
Một hình tam giác có độ dài đáy $8$ cm và có diện tích bằng diện tích hình vuông cạnh $8$ cm. Tính chiều cao của hình tam giác đó.
\end{probbox}
\probsol{
Hình vuông: $8\times8=64$ cm$^2$.\par Chiều cao: $64\times2:8=16$ cm.
\dapso{$16$ cm.}}
\gap
\begin{probbox}{29}{\relax}
Một cánh đồng hình tam giác có độ dài đáy $120$ m và chiều cao $80$ m. Tính diện tích cánh đồng đó theo đơn vị héc-ta.
\end{probbox}
\probsol{
Diện tích: $120\times80:2=4\,800$ m$^2$.\par $1$ ha $=10\,000$ m$^2$, nên diện tích đó là $0{,}48$ ha.
\dapso{$0{,}48$ ha.}}
\gap
\begin{probbox}{30}{\relax}
Hình tam giác ABC có diện tích $48$ cm$^2$; M là trung điểm của BC. Tính diện tích hình tam giác ABM.
\end{probbox}
\probsol{
ABM và ABC chung chiều cao hạ từ A.\par $BM=\dfrac12BC$, nên $S_{ABM}=\dfrac12S_{ABC}$.\par $S_{ABM}=48:2=24$ cm$^2$.
\dapso{$24$ cm$^2$.}}
\dividerbar{Thử thách -- Nâng cao}
\begin{probbox}{31}{\stars{1}}
Một hình tam giác có diện tích $120$ cm$^2$. Nếu giảm độ dài đáy đi $3$ cm (giữ nguyên chiều cao) thì diện tích giảm đi $18$ cm$^2$. Tính độ dài đáy.
\end{probbox}
\probsol{
Chiều cao: $18\times2:3=12$ cm.\par Độ dài đáy: $120\times2:12=20$ cm.
\dapso{$20$ cm.}}
\gap
\begin{probbox}{32}{\stars{1}}
Tổng độ dài hai cạnh góc vuông của một hình tam giác vuông là $14$ cm, cạnh này dài hơn cạnh kia $2$ cm. Tính diện tích hình tam giác đó.
\end{probbox}
\probsol{
Hai cạnh góc vuông dài $(14+2):2=8$ cm và $6$ cm.\par Diện tích: $8\times6:2=24$ cm$^2$.
\dapso{$24$ cm$^2$.}}
\gap
\begin{probbox}{33}{\stars{1}}
Hình chữ nhật ABCD có $AB=16$ cm và $AD=10$ cm. Điểm E nằm trên AB sao cho $AE=6$ cm. Tính diện tích các hình tam giác ADE, EBC và ECD.
\end{probbox}
\probsol{
$S_{ADE}=6\times10:2=30$ cm$^2$;\par $EB=10$ cm, nên $S_{EBC}=10\times10:2=50$ cm$^2$.\par
$S_{ECD}=16\times10:2=80$ cm$^2$ (bằng một nửa diện tích hình chữ nhật).\par \emph{Thử lại:} $30+50+80=160=16\times10$.
\dapso{$30$ cm$^2$; $50$ cm$^2$; $80$ cm$^2$.}}
\gap
\begin{probbox}{34}{\stars{2}}
Hình tam giác ABC vuông tại A có $AB=6$ cm và $AC=8$ cm. M nằm trên AB sao cho $AM=2$ cm, N nằm trên AC sao cho $AN=4$ cm. Tính diện tích hình tứ giác BMNC.
\end{probbox}
\probsol{
$S_{ABC}=6\times8:2=24$ cm$^2$;\par $S_{AMN}=2\times4:2=4$ cm$^2$.\par $S_{BMNC}=24-4=20$ cm$^2$.
\dapso{$20$ cm$^2$.}}
\gap
\begin{probbox}{35}{\stars{2}}
\begin{minipage}[t]{0.62\linewidth}
Trong hình bên có bao nhiêu hình tam giác?
\end{minipage}\hfill
\begin{minipage}[t]{0.34\linewidth}\vspace{0pt}\centering
\begin{tikzpicture}[scale=0.6]
\coordinate (A) at (2,3); \coordinate (B) at (0,0); \coordinate (C) at (4,0);
\draw[geo] (A)--(B)--(C)--cycle;
\foreach \x in {1,2,3}{\draw[geo] (A)--(\x,0);}
\draw[geo] (1,1.5)--(3,1.5);
\end{tikzpicture}
\end{minipage}
\end{probbox}
\probsol{
Chưa tính đoạn ngang ở giữa: trên đáy có $5$ điểm, tạo thành $4+3+2+1=10$ hình tam giác có đỉnh A.\par
Đoạn ngang ở giữa cũng có $5$ điểm, tạo thêm $10$ hình tam giác có đỉnh A.\par Tất cả có: $20$ hình tam giác.
\dapso{$20$ hình tam giác.}}
\gap
\begin{probbox}{36}{\stars{1}}
Hình tam giác ABC có $BC=20$ cm. M nằm trên BC sao cho $BM=8$ cm, và hình tam giác ABM có diện tích $40$ cm$^2$. Tính diện tích hình tam giác ABC.
\end{probbox}
\probsol{
Chiều cao chung hạ từ A: $40\times2:8=10$ cm.\par $S_{ABC}=20\times10:2=100$ cm$^2$.
\dapso{$100$ cm$^2$.}}
\gap
\begin{probbox}{37}{\stars{1}}
Hình chữ nhật ABCD có $AB=20$ cm, $AD=12$ cm. M là trung điểm của AB, N là trung điểm của BC. Tính diện tích hình tam giác DMN.
\end{probbox}
\probsol{
Hình chữ nhật: $20\times12=240$ cm$^2$.\par $S_{AMD}=60$, $S_{MBN}=30$, $S_{NCD}=60$ cm$^2$.\par $S_{DMN}=240-60-30-60=90$ cm$^2$.
\dapso{$90$ cm$^2$.}}
\gap
\begin{probbox}{38}{\stars{1}}
Hình vuông ABCD có cạnh $12$ cm, E là một điểm bất kì trên AB. Tính diện tích hình tam giác ECD và tổng diện tích hai hình tam giác AED và EBC.
\end{probbox}
\probsol{
Hình tam giác ECD có đáy $CD=12$ cm và chiều cao $12$ cm.\par $S_{ECD}=72$ cm$^2$.\par Tổng diện tích hai hình tam giác còn lại là $144-72=72$ cm$^2$, dù E ở vị trí nào.
\dapso{Đều bằng $72$ cm$^2$.}}
\gap
\begin{probbox}{39}{\stars{1}}
Một hình tam giác có diện tích $40$ cm$^2$. Nếu gấp độ dài đáy lên ba lần và giảm chiều cao đi một nửa thì diện tích mới là bao nhiêu?
\end{probbox}
\probsol{
Gấp đáy lên ba lần thì diện tích gấp lên ba lần.\par Giảm chiều cao đi một nửa thì diện tích giảm đi một nửa.\par Diện tích mới: $40\times3:2=60$ cm$^2$.
\dapso{$60$ cm$^2$.}}
\gap
\begin{probbox}{40}{\stars{1}}
Tổng độ dài hai cạnh góc vuông của một hình tam giác vuông là $23$ cm, hiệu của chúng là $7$ cm. Tính diện tích hình tam giác đó.
\end{probbox}
\probsol{
Cạnh góc vuông dài hơn: $(23+7):2=15$ cm.\par Cạnh ngắn hơn: $8$ cm.\par Diện tích: $15\times8:2=60$ cm$^2$.
\dapso{$60$ cm$^2$.}}
\gap
\begin{probbox}{41}{\stars{1}}
Hình tam giác ABC có diện tích $90$ cm$^2$. M nằm trên BC sao cho $MC$ gấp đôi $BM$. Tính diện tích các hình tam giác ABM và AMC.
\end{probbox}
\probsol{
$BM=\dfrac13BC$ và $MC=\dfrac23BC$; hai hình tam giác chung chiều cao hạ từ A.\par $S_{ABM}=90:3=30$ cm$^2$.\par $S_{AMC}=60$ cm$^2$.
\dapso{$S_{ABM}=30$ cm$^2$, $S_{AMC}=60$ cm$^2$.}}
\gap
\begin{probbox}{42}{\stars{1}}
Hình tam giác ABC vuông tại A có $AB=12$ cm, $AC=16$ cm. M là trung điểm của AC. So sánh diện tích hai hình tam giác ABM và BMC.
\end{probbox}
\probsol{
$AM=MC=8$ cm; cả hai hình tam giác đều có chiều cao $AB=12$ cm.\par $S_{ABM}=48$ cm$^2$.\par $S_{BMC}=48$ cm$^2$.\par Vậy diện tích hai hình tam giác bằng nhau.
\dapso{Bằng nhau, mỗi hình $48$ cm$^2$.}}
\gap
\begin{probbox}{43}{\stars{1}}
Hình tam giác ABC có $BC=24$ cm và chiều cao $AH=10$ cm. Kéo dài BC về phía C một đoạn $CD=6$ cm. Tính diện tích các hình tam giác ACD và ABD.
\end{probbox}
\probsol{
Các hình tam giác ABC, ACD và ABD đều có chiều cao $AH=10$ cm.\par $S_{ACD}=6\times10:2=30$ cm$^2$.\par $S_{ABD}=30\times10:2=150$ cm$^2$.
\dapso{$30$ cm$^2$ và $150$ cm$^2$.}}
\gap
\begin{probbox}{44}{\stars{1}}
Hai hình tam giác có chiều cao bằng nhau. Đáy của hình thứ nhất bằng $\dfrac23$ đáy của hình thứ hai, và diện tích hình thứ hai lớn hơn diện tích hình thứ nhất $24$ cm$^2$. Tính diện tích mỗi hình.
\end{probbox}
\probsol{
Khi chiều cao bằng nhau, tỉ số diện tích bằng tỉ số độ dài đáy: $S_1=\dfrac23S_2$.\par Hiệu số phần bằng nhau là $3-2=1$ phần, ứng với $24$ cm$^2$.\par $S_1=48$ cm$^2$; $S_2=72$ cm$^2$.
\dapso{$48$ cm$^2$ và $72$ cm$^2$.}}
\gap
\begin{probbox}{45}{\stars{1}}
Một mảnh vườn hình chữ nhật có chiều dài $40$ m, chiều rộng $25$ m. Người ta trồng hoa trên một phần đất hình tam giác có đáy là một chiều dài của mảnh vườn và đỉnh nằm trên cạnh đối diện; phần còn lại trồng cỏ. Tính diện tích mỗi phần.
\end{probbox}
\probsol{
Hình tam giác có đáy $40$ m và chiều cao $25$ m.\par Phần trồng hoa: $500$ m$^2$.\par Mảnh vườn: $1\,000$ m$^2$.\par Phần trồng cỏ: $500$ m$^2$.
\dapso{Trồng hoa $500$ m$^2$, trồng cỏ $500$ m$^2$.}}
\gap
\begin{probbox}{46}{\stars{1}}
Hình tam giác ABC có diện tích $72$ cm$^2$. D nằm trên AB sao cho $AD=\dfrac13AB$; E là trung điểm của AC. Tính diện tích hình tam giác ADE.
\end{probbox}
\probsol{
Nối D với C.\par ADC và ABC chung chiều cao hạ từ C, và $AD=\dfrac13AB$.\par $S_{ADC}=72:3=24$ cm$^2$.\par ADE và ADC chung chiều cao hạ từ D, và $AE=\dfrac12AC$.\par $S_{ADE}=12$ cm$^2$.
\dapso{$12$ cm$^2$.}}
\gap
\begin{probbox}{47}{\stars{1}}
Hình vuông ABCD có cạnh $10$ cm. E là trung điểm của BC, F là trung điểm của CD. Tính diện tích hình tam giác AEF.
\end{probbox}
\probsol{
$S_{ABE}=25$ cm$^2$; $S_{ADF}=25$ cm$^2$.\par $S_{ECF}=5\times5:2=12{,}5$ cm$^2$.\par $S_{AEF}=100-25-25-12{,}5=37{,}5$ cm$^2$.
\dapso{$37{,}5$ cm$^2$.}}
\gap
\begin{probbox}{48}{\stars{2}}
Một hình vuông có đường chéo dài $8$ cm. Tính diện tích hình vuông đó.
\end{probbox}
\probsol{
Một đường chéo chia hình vuông thành hai hình tam giác bằng nhau có chung đáy $8$ cm.\par Hai đường chéo bằng nhau, vuông góc với nhau và cắt nhau tại trung điểm của mỗi đường, nên chiều cao của mỗi hình tam giác là $8:2=4$ cm.\par Mỗi hình tam giác: $8\times4:2=16$ cm$^2$.\par Hình vuông: $32$ cm$^2$.
\dapso{$32$ cm$^2$.}}
\gap
\begin{probbox}{49}{\stars{2}}
Hình tam giác ABC có diện tích $80$ cm$^2$. M là trung điểm của BC, N là trung điểm của AM. Tính diện tích hình tam giác BNC.
\end{probbox}
\probsol{
$S_{ABM}=S_{ACM}=40$ cm$^2$.\par BNM và BAM chung chiều cao hạ từ B và $NM=\dfrac12AM$, nên $S_{BNM}=20$ cm$^2$.\par Tương tự, $S_{CNM}=20$ cm$^2$.\par $S_{BNC}=40$ cm$^2$.
\dapso{$40$ cm$^2$.}}
\gap
\begin{probbox}{50}{\stars{2}}
Hình tam giác ABC có diện tích $120$ cm$^2$. M nằm trên BC sao cho $BM=\dfrac14BC$; N nằm trên AM sao cho $AN=\dfrac23AM$. Tính diện tích các hình tam giác ABN và NBM.
\end{probbox}
\probsol{
$S_{ABM}=120:4=30$ cm$^2$ (chung chiều cao hạ từ A).\par ABN và ABM chung chiều cao hạ từ B, với $AN=\dfrac23AM$.\par $S_{ABN}=20$ cm$^2$.\par $S_{NBM}=10$ cm$^2$.
\dapso{$S_{ABN}=20$ cm$^2$, $S_{NBM}=10$ cm$^2$.}}
\clearpage
\titlehead
\begingroup\setbox0=\vbox{\section*{(*) So sánh diện tích hình tam giác}}\endgroup
\chapterhead{12}{(*) So sánh diện tích hình tam giác}

\shead{Chung chiều cao, so sánh đáy}
\begin{knowbox}
\begin{minipage}[t]{0.6\linewidth}
\begin{itemize}[leftmargin=1.4em,itemsep=3pt]
\item Hai hình tam giác có \textbf{chung chiều cao} thì tỉ số diện tích bằng tỉ số \textbf{độ dài hai đáy}. Chẳng hạn, nếu M nằm trên BC thì
\[\frac{S_{ABM}}{S_{AMC}}=\frac{BM}{MC}.\]
\item Hai hình tam giác có \textbf{chung đáy} thì tỉ số diện tích bằng tỉ số \textbf{hai chiều cao}.
\item Đoạn thẳng nối một đỉnh với \textbf{trung điểm} của cạnh đối diện chia hình tam giác thành hai hình tam giác có diện tích bằng nhau.
\end{itemize}
Để so sánh hai hình tam giác, ta thường đi qua một hình tam giác ``trung gian'' có chung chiều cao với từng hình.
\end{minipage}\hfill
\begin{minipage}[t]{0.36\linewidth}\vspace{4pt}\centering
\begin{tikzpicture}[scale=0.85]
\coordinate (B) at (0,0); \coordinate (C) at (4.2,0); \coordinate (A) at (1.4,2.6); \coordinate (M) at (2.8,0); \coordinate (H) at (1.4,0);
\fill[tmBlue!18] (A)--(B)--(M)--cycle; \fill[tmOrange!18] (A)--(M)--(C)--cycle;
\draw[geo] (A)--(B)--(C)--cycle; \draw[geo] (A)--(M);
\draw[dashed,Slate] (A)--(H); \ramark{H}{C}{A}
\node[lbl,above] at (A){A}; \node[lbl,below] at (B){B}; \node[lbl,below] at (C){C}; \node[lbl,below] at (M){M};
\end{tikzpicture}
\end{minipage}
\end{knowbox}
\gap
\begin{exbox}{\ExWord\ 1}
Hình tam giác ABC có diện tích $45$ cm$^2$. Điểm M nằm trên BC sao cho $BM$ gấp đôi $MC$. Tính diện tích các hình tam giác ABM và AMC.
\guide
Hai hình tam giác ABM và AMC chung chiều cao hạ từ A, nên tỉ số diện tích của chúng là $BM:MC=2:1$.\par Nếu $S_{AMC}$ là $1$ phần thì $S_{ABM}$ là $2$ phần và $S_{ABC}$ là $3$ phần.\par
$S_{AMC}=45:3=15$ cm$^2$;\par $S_{ABM}=15\times2=30$ cm$^2$.
\dapso{$30$ cm$^2$ và $15$ cm$^2$.}
\end{exbox}
\gap
\begin{exbox}{\ExWord\ 2}
\begin{minipage}[t]{0.6\linewidth}
Hình tam giác ABC có diện tích $60$ cm$^2$; M là trung điểm của BC và N là trung điểm của AM. Tính diện tích các hình tam giác ABN và BNC.
\end{minipage}\hfill
\begin{minipage}[t]{0.36\linewidth}\vspace{0pt}\centering
\begin{tikzpicture}[scale=0.7]
\coordinate (B) at (0,0); \coordinate (C) at (4,0); \coordinate (A) at (1.2,2.6); \coordinate (M) at (2,0); \coordinate (N) at (1.6,1.3);
\fill[shade] (B)--(N)--(C)--cycle;
\draw[geo] (A)--(B)--(C)--cycle; \draw[geo] (A)--(M) (B)--(N)--(C);
\node[lbl,above] at (A){A}; \node[lbl,below] at (B){B}; \node[lbl,below] at (C){C}; \node[lbl,below] at (M){M}; \node[lbl,right] at (N){N};
\end{tikzpicture}
\end{minipage}
\guide
M là trung điểm của BC, nên $S_{ABM}=S_{AMC}=60:2=30$ cm$^2$.\par
N là trung điểm của AM, nên $S_{ABN}=S_{NBM}=30:2=15$ cm$^2$; tương tự, $S_{NMC}=15$ cm$^2$.\par
$S_{BNC}=S_{NBM}+S_{NMC}=15+15=30$ cm$^2$.
\dapso{$S_{ABN}=15$ cm$^2$;\par $S_{BNC}=30$ cm$^2$.}
\end{exbox}
\gap
\begin{exbox}{\ExWord\ 3}
Hình tam giác ABC có diện tích $72$ cm$^2$. M nằm trên AB sao cho $AM=\dfrac13AB$, N nằm trên AC sao cho $AN=\dfrac12AC$. Tính diện tích hình tam giác AMN và hình tứ giác BMNC.
\guide
Nối M với C (hình tam giác ``trung gian'' AMC).\par
AMC và ABC chung chiều cao hạ từ C, và $AM=\dfrac13AB$, nên $S_{AMC}=\dfrac13\times72=24$ cm$^2$.\par
AMN và AMC chung chiều cao hạ từ M, và $AN=\dfrac12AC$, nên $S_{AMN}=\dfrac12\times24=12$ cm$^2$.\par
$S_{BMNC}=72-12=60$ cm$^2$.
\dapso{$12$ cm$^2$; $60$ cm$^2$.}
\end{exbox}
\gap
\begin{exbox}{\ExWord\ 4}
Hình vuông ABCD có cạnh $12$ cm; M là trung điểm của AB, N là trung điểm của BC. Tính diện tích hình tam giác DMN.
\guide
Lấy diện tích hình vuông trừ đi diện tích ba hình tam giác vuông ở các góc:\par
$S_{AMD}=6\times12:2=36$;\par $S_{MBN}=6\times6:2=18$;\par $S_{NCD}=6\times12:2=36$ cm$^2$.\par
$S_{DMN}=12\times12-36-18-36=54$ cm$^2$.
\dapso{$54$ cm$^2$.}
\end{exbox}
\dividerbar{Luyện tập}
\begin{probbox}{1}{\relax}
Hình tam giác ABC có diện tích $40$ cm$^2$. M nằm trên BC sao cho $BM=3$ cm và $MC=5$ cm. Tính diện tích các hình tam giác ABM và AMC.
\end{probbox}
\probsol{
Tỉ số diện tích là $3:5$, tổng cộng $8$ phần.\par $S_{ABM}=40:8\times3=15$ cm$^2$; $S_{AMC}=25$ cm$^2$.
\dapso{$15$ cm$^2$ và $25$ cm$^2$.}}
\gap
\begin{probbox}{2}{\relax}
Trong hình tam giác ABC, M là trung điểm của BC và $S_{ABM}=18$ cm$^2$. Tính $S_{ABC}$.
\end{probbox}
\probsol{
$S_{AMC}=S_{ABM}=18$ cm$^2$, nên $S_{ABC}=36$ cm$^2$.
\dapso{$36$ cm$^2$.}}
\gap
\begin{probbox}{3}{\relax}
Hình tam giác ABC có diện tích $90$ cm$^2$. D nằm trên AC sao cho $AD=\dfrac23AC$. Tính diện tích các hình tam giác ABD và BDC.
\end{probbox}
\probsol{
Hai hình tam giác chung chiều cao hạ từ B: $S_{ABD}=\dfrac23\times90=60$ cm$^2$ và $S_{BDC}=30$ cm$^2$.
\dapso{$60$ cm$^2$ và $30$ cm$^2$.}}
\gap
\begin{probbox}{4}{\relax}
Hai điểm M và N chia cạnh BC của hình tam giác ABC thành ba phần bằng nhau $BM=MN=NC$. Biết $S_{AMN}=15$ cm$^2$, tính $S_{ABC}$.
\end{probbox}
\probsol{
ABM, AMN và ANC chung chiều cao và có đáy bằng nhau, nên diện tích của chúng bằng nhau.\par $S_{ABC}=15\times3=45$ cm$^2$.
\dapso{$45$ cm$^2$.}}
\gap
\begin{probbox}{5}{\relax}
Hình tam giác ABC có diện tích $20$ cm$^2$. Kéo dài BC về phía C đến D sao cho $CD=BC$. Tính diện tích các hình tam giác ACD và ABD.
\end{probbox}
\probsol{
ACD và ABC chung chiều cao hạ từ A và có đáy bằng nhau: $S_{ACD}=20$ cm$^2$.\par $S_{ABD}=20+20=40$ cm$^2$.
\dapso{$20$ cm$^2$; $40$ cm$^2$.}}
\gap
\begin{probbox}{6}{\relax}
Hình tam giác ABC có diện tích $60$ cm$^2$; M là trung điểm của BC; N nằm trên AC sao cho $AN=\dfrac13AC$. Tính $S_{MNC}$.
\end{probbox}
\probsol{
$S_{AMC}=60:2=30$ cm$^2$.\par MNC và AMC chung chiều cao hạ từ M và $NC=\dfrac23AC$, nên $S_{MNC}=\dfrac23\times30=20$ cm$^2$.
\dapso{$20$ cm$^2$.}}
\gap
\begin{probbox}{7}{\relax}
Hình tam giác ABC có diện tích $54$ cm$^2$. M nằm trên BC sao cho $BM=\dfrac13BC$; N là trung điểm của AM. Tính diện tích các hình tam giác BNM và ANC.
\end{probbox}
\probsol{
$S_{ABM}=\dfrac13\times54=18$ và $S_{AMC}=36$ cm$^2$.\par N là trung điểm của AM, nên $S_{BNM}=18:2=9$ cm$^2$ và $S_{ANC}=36:2=18$ cm$^2$.
\dapso{$9$ cm$^2$; $18$ cm$^2$.}}
\gap
\begin{probbox}{8}{\relax}
Hình tam giác ABC có diện tích $48$ cm$^2$. M nằm trên AB sao cho $AM=\dfrac23AB$; N nằm trên AC sao cho $AN=\dfrac34AC$. Tính diện tích hình tam giác AMN và hình tứ giác BMNC.
\end{probbox}
\probsol{
$S_{AMC}=\dfrac23\times48=32$ cm$^2$;\par $S_{AMN}=\dfrac34\times32=24$ cm$^2$;\par $S_{BMNC}=48-24=24$ cm$^2$.
\dapso{$24$ cm$^2$; $24$ cm$^2$.}}
\gap
\begin{probbox}{9}{\relax}
Trong hình tam giác ABC, D nằm trên BC sao cho $BD=\dfrac14BC$, và $S_{ABD}=12$ cm$^2$. Tính $S_{ADC}$.
\end{probbox}
\probsol{
$DC=3\times BD$ và hai hình tam giác chung chiều cao hạ từ A, nên $S_{ADC}=12\times3=36$ cm$^2$.
\dapso{$36$ cm$^2$.}}
\gap
\begin{probbox}{10}{\relax}
Hình tam giác ABC có diện tích $60$ cm$^2$. Điểm M trên BC chia BC sao cho $BM:MC=1:2$. Tính diện tích các hình tam giác ABM và AMC.
\end{probbox}
\probsol{
ABM, AMC và ABC chung chiều cao hạ từ A.\par $BM=\dfrac{1}{3}BC$, nên $S_{ABM}=60\times\dfrac{1}{3}=20$ cm$^2$.\par $S_{AMC}=60-20=40$ cm$^2$.
\dapso{$S_{ABM}=20$ cm$^2$, $S_{AMC}=40$ cm$^2$.}}
\gap
\begin{probbox}{11}{\relax}
Hình tam giác ABC có diện tích $84$ cm$^2$. Điểm M trên BC chia BC sao cho $BM:MC=3:4$. Tính diện tích các hình tam giác ABM và AMC.
\end{probbox}
\probsol{
ABM, AMC và ABC chung chiều cao hạ từ A.\par $BM=\dfrac{3}{7}BC$, nên $S_{ABM}=84\times\dfrac{3}{7}=36$ cm$^2$.\par $S_{AMC}=84-36=48$ cm$^2$.
\dapso{$S_{ABM}=36$ cm$^2$, $S_{AMC}=48$ cm$^2$.}}
\gap
\begin{probbox}{12}{\relax}
Hình tam giác ABC có diện tích $100$ cm$^2$. Điểm M trên BC chia BC sao cho $BM:MC=2:3$. Tính diện tích các hình tam giác ABM và AMC.
\end{probbox}
\probsol{
ABM, AMC và ABC chung chiều cao hạ từ A.\par $BM=\dfrac{2}{5}BC$, nên $S_{ABM}=100\times\dfrac{2}{5}=40$ cm$^2$.\par $S_{AMC}=100-40=60$ cm$^2$.
\dapso{$S_{ABM}=40$ cm$^2$, $S_{AMC}=60$ cm$^2$.}}
\gap
\begin{probbox}{13}{\relax}
Hình tam giác ABC có diện tích $72$ cm$^2$. Điểm M trên BC chia BC sao cho $BM:MC=5:1$. Tính diện tích các hình tam giác ABM và AMC.
\end{probbox}
\probsol{
ABM, AMC và ABC chung chiều cao hạ từ A.\par $BM=\dfrac{5}{6}BC$, nên $S_{ABM}=72\times\dfrac{5}{6}=60$ cm$^2$.\par $S_{AMC}=72-60=12$ cm$^2$.
\dapso{$S_{ABM}=60$ cm$^2$, $S_{AMC}=12$ cm$^2$.}}
\gap
\begin{probbox}{14}{\relax}
Hình tam giác ABC có diện tích $45$ cm$^2$. Điểm M trên BC chia BC sao cho $BM:MC=4:5$. Tính diện tích các hình tam giác ABM và AMC.
\end{probbox}
\probsol{
ABM, AMC và ABC chung chiều cao hạ từ A.\par $BM=\dfrac{4}{9}BC$, nên $S_{ABM}=45\times\dfrac{4}{9}=20$ cm$^2$.\par $S_{AMC}=45-20=25$ cm$^2$.
\dapso{$S_{ABM}=20$ cm$^2$, $S_{AMC}=25$ cm$^2$.}}
\gap
\begin{probbox}{15}{\relax}
Trong hình tam giác ABC, điểm M trên BC có $BM=\dfrac{1}{3}BC$. Hình tam giác ABM có diện tích $15$ cm$^2$. Tính diện tích các hình tam giác ABC và AMC.
\end{probbox}
\probsol{
ABM và ABC chung chiều cao hạ từ A, nên $S_{ABM}=\dfrac{1}{3}S_{ABC}$.\par $S_{ABC}=15:\dfrac{1}{3}=45$ cm$^2$.\par $S_{AMC}=45-15=30$ cm$^2$.
\dapso{$S_{ABC}=45$ cm$^2$, $S_{AMC}=30$ cm$^2$.}}
\gap
\begin{probbox}{16}{\relax}
Trong hình tam giác ABC, điểm M trên BC có $BM=\dfrac{2}{5}BC$. Hình tam giác ABM có diện tích $24$ cm$^2$. Tính diện tích các hình tam giác ABC và AMC.
\end{probbox}
\probsol{
ABM và ABC chung chiều cao hạ từ A, nên $S_{ABM}=\dfrac{2}{5}S_{ABC}$.\par $S_{ABC}=24:\dfrac{2}{5}=60$ cm$^2$.\par $S_{AMC}=60-24=36$ cm$^2$.
\dapso{$S_{ABC}=60$ cm$^2$, $S_{AMC}=36$ cm$^2$.}}
\gap
\begin{probbox}{17}{\relax}
Trong hình tam giác ABC, điểm M trên BC có $BM=\dfrac{3}{4}BC$. Hình tam giác ABM có diện tích $21$ cm$^2$. Tính diện tích các hình tam giác ABC và AMC.
\end{probbox}
\probsol{
ABM và ABC chung chiều cao hạ từ A, nên $S_{ABM}=\dfrac{3}{4}S_{ABC}$.\par $S_{ABC}=21:\dfrac{3}{4}=28$ cm$^2$.\par $S_{AMC}=28-21=7$ cm$^2$.
\dapso{$S_{ABC}=28$ cm$^2$, $S_{AMC}=7$ cm$^2$.}}
\gap
\begin{probbox}{18}{\relax}
Trong hình tam giác ABC, điểm M trên BC có $BM=\dfrac{1}{6}BC$. Hình tam giác ABM có diện tích $12$ cm$^2$. Tính diện tích các hình tam giác ABC và AMC.
\end{probbox}
\probsol{
ABM và ABC chung chiều cao hạ từ A, nên $S_{ABM}=\dfrac{1}{6}S_{ABC}$.\par $S_{ABC}=12:\dfrac{1}{6}=72$ cm$^2$.\par $S_{AMC}=72-12=60$ cm$^2$.
\dapso{$S_{ABC}=72$ cm$^2$, $S_{AMC}=60$ cm$^2$.}}
\gap
\begin{probbox}{19}{\relax}
Hình tam giác ABC có diện tích $80$ cm$^2$ và $BC=16$ cm. Điểm M trên BC sao cho hình tam giác ABM có diện tích $30$ cm$^2$. Tính BM.
\end{probbox}
\probsol{
ABM và ABC chung chiều cao hạ từ A, nên tỉ số diện tích bằng tỉ số độ dài hai đáy.\par $\dfrac{BM}{BC}=\dfrac{30}{80}=\dfrac{3}{8}$.\par $BM=16\times\dfrac{3}{8}=6$ cm.
\dapso{$BM=6$ cm.}}
\gap
\begin{probbox}{20}{\relax}
Hình tam giác ABC có diện tích $96$ cm$^2$ và $BC=12$ cm. Điểm M trên BC sao cho hình tam giác ABM có diện tích $36$ cm$^2$. Tính BM.
\end{probbox}
\probsol{
ABM và ABC chung chiều cao hạ từ A, nên tỉ số diện tích bằng tỉ số độ dài hai đáy.\par $\dfrac{BM}{BC}=\dfrac{36}{96}=\dfrac{3}{8}$.\par $BM=12\times\dfrac{3}{8}=4{,}5$ cm.
\dapso{$BM=4{,}5$ cm.}}
\gap
\begin{probbox}{21}{\relax}
Hình tam giác ABC có diện tích $75$ cm$^2$ và $BC=10$ cm. Điểm M trên BC sao cho hình tam giác ABM có diện tích $45$ cm$^2$. Tính BM.
\end{probbox}
\probsol{
ABM và ABC chung chiều cao hạ từ A, nên tỉ số diện tích bằng tỉ số độ dài hai đáy.\par $\dfrac{BM}{BC}=\dfrac{45}{75}=\dfrac{3}{5}$.\par $BM=10\times\dfrac{3}{5}=6$ cm.
\dapso{$BM=6$ cm.}}
\gap
\begin{probbox}{22}{\relax}
Hình tam giác ABC có diện tích $36$ cm$^2$. Kéo dài BC về phía C một đoạn $CD=\dfrac12BC$. Tính diện tích các hình tam giác ACD và ABD.
\end{probbox}
\probsol{
ABC, ACD và ABD chung chiều cao hạ từ A.\par $S_{ACD}=36\times\dfrac{1}{2}=18$ cm$^2$.\par $S_{ABD}=36+18=54$ cm$^2$.
\dapso{$S_{ACD}=18$ cm$^2$, $S_{ABD}=54$ cm$^2$.}}
\gap
\begin{probbox}{23}{\relax}
Hình tam giác ABC có diện tích $50$ cm$^2$. Kéo dài BC về phía C một đoạn $CD=\dfrac25BC$. Tính diện tích các hình tam giác ACD và ABD.
\end{probbox}
\probsol{
ABC, ACD và ABD chung chiều cao hạ từ A.\par $S_{ACD}=50\times\dfrac{2}{5}=20$ cm$^2$.\par $S_{ABD}=50+20=70$ cm$^2$.
\dapso{$S_{ACD}=20$ cm$^2$, $S_{ABD}=70$ cm$^2$.}}
\gap
\begin{probbox}{24}{\relax}
Hình tam giác ABC có diện tích $28$ cm$^2$. Kéo dài BC về phía C một đoạn $CD=2\times BC$. Tính diện tích các hình tam giác ACD và ABD.
\end{probbox}
\probsol{
ABC, ACD và ABD chung chiều cao hạ từ A.\par $S_{ACD}=28\times2=56$ cm$^2$.\par $S_{ABD}=28+56=84$ cm$^2$.
\dapso{$S_{ACD}=56$ cm$^2$, $S_{ABD}=84$ cm$^2$.}}
\gap
\begin{probbox}{25}{\relax}
Hình tam giác ABC có diện tích $60$ cm$^2$. Điểm M trên AB có $AM=\dfrac{1}{2}AB$; điểm N trên AC có $AN=\dfrac{1}{3}AC$. Tính diện tích hình tam giác AMN và hình tứ giác BMNC.
\end{probbox}
\probsol{
Nối M với C.\par AMC và ABC chung chiều cao hạ từ C: $S_{AMC}=60\times\dfrac{1}{2}=30$ cm$^2$.\par AMN và AMC chung chiều cao hạ từ M: $S_{AMN}=30\times\dfrac{1}{3}=10$ cm$^2$.\par $S_{BMNC}=60-10=50$ cm$^2$.
\dapso{$S_{AMN}=10$ cm$^2$, $S_{BMNC}=50$ cm$^2$.}}
\gap
\begin{probbox}{26}{\relax}
Hình tam giác ABC có diện tích $90$ cm$^2$. Điểm M trên AB có $AM=\dfrac{2}{3}AB$; điểm N trên AC có $AN=\dfrac{1}{2}AC$. Tính diện tích hình tam giác AMN và hình tứ giác BMNC.
\end{probbox}
\probsol{
Nối M với C.\par AMC và ABC chung chiều cao hạ từ C: $S_{AMC}=90\times\dfrac{2}{3}=60$ cm$^2$.\par AMN và AMC chung chiều cao hạ từ M: $S_{AMN}=60\times\dfrac{1}{2}=30$ cm$^2$.\par $S_{BMNC}=90-30=60$ cm$^2$.
\dapso{$S_{AMN}=30$ cm$^2$, $S_{BMNC}=60$ cm$^2$.}}
\gap
\begin{probbox}{27}{\relax}
Hình tam giác ABC có diện tích $84$ cm$^2$. Điểm M trên AB có $AM=\dfrac{3}{4}AB$; điểm N trên AC có $AN=\dfrac{2}{3}AC$. Tính diện tích hình tam giác AMN và hình tứ giác BMNC.
\end{probbox}
\probsol{
Nối M với C.\par AMC và ABC chung chiều cao hạ từ C: $S_{AMC}=84\times\dfrac{3}{4}=63$ cm$^2$.\par AMN và AMC chung chiều cao hạ từ M: $S_{AMN}=63\times\dfrac{2}{3}=42$ cm$^2$.\par $S_{BMNC}=84-42=42$ cm$^2$.
\dapso{$S_{AMN}=42$ cm$^2$, $S_{BMNC}=42$ cm$^2$.}}
\gap
\begin{probbox}{28}{\relax}
Hai hình tam giác ABC và DBC có chung đáy BC. Chiều cao hạ từ A là $8$ cm, chiều cao hạ từ D là $6$ cm. Biết $S_{ABC}=40$ cm$^2$, tính $S_{DBC}$.
\end{probbox}
\probsol{
Khi chung đáy, tỉ số diện tích bằng tỉ số hai chiều cao: $\dfrac{S_{DBC}}{S_{ABC}}=\dfrac{6}{8}$.\par $S_{DBC}=40\times6:8=30$ cm$^2$.
\dapso{$30$ cm$^2$.}}
\gap
\begin{probbox}{29}{\relax}
Hai hình tam giác ABC và DBC có chung đáy BC. Chiều cao hạ từ A là $9$ cm, chiều cao hạ từ D là $12$ cm. Biết $S_{ABC}=54$ cm$^2$, tính $S_{DBC}$.
\end{probbox}
\probsol{
Khi chung đáy, tỉ số diện tích bằng tỉ số hai chiều cao: $\dfrac{S_{DBC}}{S_{ABC}}=\dfrac{12}{9}$.\par $S_{DBC}=54\times12:9=72$ cm$^2$.
\dapso{$72$ cm$^2$.}}
\gap
\begin{probbox}{30}{\relax}
Hai hình tam giác ABC và DBC có chung đáy BC. Chiều cao hạ từ A là $10$ cm, chiều cao hạ từ D là $4$ cm. Biết $S_{ABC}=35$ cm$^2$, tính $S_{DBC}$.
\end{probbox}
\probsol{
Khi chung đáy, tỉ số diện tích bằng tỉ số hai chiều cao: $\dfrac{S_{DBC}}{S_{ABC}}=\dfrac{4}{10}$.\par $S_{DBC}=35\times4:10=14$ cm$^2$.
\dapso{$14$ cm$^2$.}}
\dividerbar{Thử thách -- Nâng cao}
\begin{probbox}{31}{\stars{1}}
Hình tam giác ABC có diện tích $64$ cm$^2$; M và N lần lượt là trung điểm của AB và AC. Tính diện tích hình tam giác AMN và hình tứ giác BMNC.
\end{probbox}
\probsol{
$S_{AMC}=64:2=32$ cm$^2$; $S_{AMN}=32:2=16$ cm$^2$.\par $S_{BMNC}=64-16=48$ cm$^2$.
\note{Diện tích hình tam giác AMN bằng $\dfrac14$ diện tích hình tam giác ABC.}
\dapso{$16$ cm$^2$; $48$ cm$^2$.}}
\gap
\begin{probbox}{32}{\stars{1}}
Hình tam giác ABC có diện tích $30$ cm$^2$. Kéo dài AB về phía B đến D sao cho $BD=AB$, kéo dài AC về phía C đến E sao cho $CE=AC$. Tính $S_{ADE}$.
\end{probbox}
\probsol{
$AD=2\times AB$, nên $S_{ADC}=2\times30=60$ cm$^2$.\par $AE=2\times AC$, nên $S_{ADE}=2\times60=120$ cm$^2$.
\dapso{$120$ cm$^2$.}}
\gap
\begin{probbox}{33}{\stars{1}}
Hình tam giác ABC có diện tích $120$ cm$^2$. D nằm trên BC sao cho $BD=2\times DC$; E là trung điểm của AD. Tính $S_{BEC}$.
\end{probbox}
\probsol{
$S_{ABD}=80$ và $S_{ADC}=40$ cm$^2$.\par E là trung điểm của AD, nên $S_{BED}=40$ và $S_{EDC}=20$ cm$^2$.\par
$S_{BEC}=40+20=60$ cm$^2$.
\dapso{$60$ cm$^2$.}}
\gap
\begin{probbox}{34}{\stars{2}}
\begin{minipage}[t]{0.6\linewidth}
Cho hình tam giác ABC, M là trung điểm của BC và I là trung điểm của AM. Đường thẳng BI cắt AC tại N. Chứng tỏ rằng $AN=\dfrac13AC$.
\end{minipage}\hfill
\begin{minipage}[t]{0.36\linewidth}\vspace{0pt}\centering
\begin{tikzpicture}[scale=0.7]
\coordinate (B) at (0,0); \coordinate (C) at (4.2,0); \coordinate (A) at (1.2,2.7); \coordinate (M) at (2.1,0); \coordinate (I) at (1.65,1.35); \coordinate (N) at (2.2,1.8);
\draw[geo] (A)--(B)--(C)--cycle; \draw[geo] (A)--(M) (B)--(N);
\draw[dashed,Slate] (M)--(N);
\node[lbl,above] at (A){A}; \node[lbl,below] at (B){B}; \node[lbl,below] at (C){C}; \node[lbl,below] at (M){M}; \node[lbl,left] at (I){I}; \node[lbl,above right] at (N){N};
\end{tikzpicture}
\end{minipage}
\end{probbox}
\probsol{
Nối M với N.\par Vì $AI=IM$ nên $S_{ABI}=S_{MBI}$ và $S_{AIN}=S_{MIN}$.\par Cộng lại: $S_{ABN}=S_{MBN}$.\par
Vì $BM=MC$ nên $S_{MBN}=S_{MNC}$.\par Do đó $S_{BNC}=2\times S_{ABN}$.\par
ABN và CBN chung chiều cao hạ từ B, nên $NC=2\times AN$, tức là\par $AN=\dfrac13AC$.
\dapso{Đã chứng tỏ được điều phải chứng minh.}}
\gap
\begin{probbox}{35}{\stars{2}}
Hình tam giác ABC có diện tích $90$ cm$^2$. M nằm trên BC sao cho $BM=\dfrac13BC$, N nằm trên CA sao cho $CN=\dfrac13CA$, và P nằm trên AB sao cho $AP=\dfrac13AB$. Tính $S_{MNP}$.
\end{probbox}
\probsol{
Hình tam giác APN: $S_{APC}=\dfrac13\times90=30$; $AN=\dfrac23AC$, nên $S_{APN}=\dfrac23\times30=20$ cm$^2$.\par
Tương tự, $S_{BMP}=20$ và $S_{CNM}=20$ cm$^2$.\par Vậy $S_{MNP}=90-3\times20=30$ cm$^2$.
\dapso{$30$ cm$^2$.}}
\gap
\begin{probbox}{36}{\stars{1}}
Hình tam giác ABC có diện tích $60$ cm$^2$. Điểm M trên BC có $BM=\dfrac{1}{3}BC$; điểm N trên AM có $AN=\dfrac{1}{2}AM$. Tính diện tích các hình tam giác ABN, NBM và ANC.
\end{probbox}
\probsol{
$S_{ABM}=60\times\dfrac{1}{3}=20$ cm$^2$; $S_{AMC}=40$ cm$^2$.\par ABN và ABM chung chiều cao hạ từ B: $S_{ABN}=20\times\dfrac{1}{2}=10$ cm$^2$.\par $S_{NBM}=20-10=10$ cm$^2$.\par ANC và AMC chung chiều cao hạ từ C: $S_{ANC}=40\times\dfrac{1}{2}=20$ cm$^2$.
\dapso{$S_{ABN}=10$, $S_{NBM}=10$, $S_{ANC}=20$ (cm$^2$).}}
\gap
\begin{probbox}{37}{\stars{1}}
Hình tam giác ABC có diện tích $96$ cm$^2$. Điểm M trên BC có $BM=\dfrac{1}{4}BC$; điểm N trên AM có $AN=\dfrac{2}{3}AM$. Tính diện tích các hình tam giác ABN, NBM và ANC.
\end{probbox}
\probsol{
$S_{ABM}=96\times\dfrac{1}{4}=24$ cm$^2$; $S_{AMC}=72$ cm$^2$.\par ABN và ABM chung chiều cao hạ từ B: $S_{ABN}=24\times\dfrac{2}{3}=16$ cm$^2$.\par $S_{NBM}=24-16=8$ cm$^2$.\par ANC và AMC chung chiều cao hạ từ C: $S_{ANC}=72\times\dfrac{2}{3}=48$ cm$^2$.
\dapso{$S_{ABN}=16$, $S_{NBM}=8$, $S_{ANC}=48$ (cm$^2$).}}
\gap
\begin{probbox}{38}{\stars{1}}
Hình tam giác ABC có diện tích $150$ cm$^2$. Điểm M trên BC có $BM=\dfrac{2}{5}BC$; điểm N trên AM có $AN=\dfrac{3}{5}AM$. Tính diện tích các hình tam giác ABN, NBM và ANC.
\end{probbox}
\probsol{
$S_{ABM}=150\times\dfrac{2}{5}=60$ cm$^2$; $S_{AMC}=90$ cm$^2$.\par ABN và ABM chung chiều cao hạ từ B: $S_{ABN}=60\times\dfrac{3}{5}=36$ cm$^2$.\par $S_{NBM}=60-36=24$ cm$^2$.\par ANC và AMC chung chiều cao hạ từ C: $S_{ANC}=90\times\dfrac{3}{5}=54$ cm$^2$.
\dapso{$S_{ABN}=36$, $S_{NBM}=24$, $S_{ANC}=54$ (cm$^2$).}}
\gap
\begin{probbox}{39}{\stars{1}}
Điểm M trên AB có $AM=\dfrac{1}{2}AB$ và điểm N trên AC có $AN=\dfrac{2}{5}AC$. Hình tam giác AMN có diện tích $10$ cm$^2$. Tính diện tích hình tam giác ABC.
\end{probbox}
\probsol{
Nối M với C.\par AMN và AMC chung chiều cao hạ từ M: $S_{AMC}=10:\dfrac{2}{5}=25$ cm$^2$.\par AMC và ABC chung chiều cao hạ từ C: $S_{ABC}=25:\dfrac{1}{2}=50$ cm$^2$.
\dapso{$50$ cm$^2$.}}
\gap
\begin{probbox}{40}{\stars{1}}
Điểm M trên AB có $AM=\dfrac{3}{4}AB$ và điểm N trên AC có $AN=\dfrac{2}{3}AC$. Hình tam giác AMN có diện tích $18$ cm$^2$. Tính diện tích hình tam giác ABC.
\end{probbox}
\probsol{
Nối M với C.\par AMN và AMC chung chiều cao hạ từ M: $S_{AMC}=18:\dfrac{2}{3}=27$ cm$^2$.\par AMC và ABC chung chiều cao hạ từ C: $S_{ABC}=27:\dfrac{3}{4}=36$ cm$^2$.
\dapso{$36$ cm$^2$.}}
\gap
\begin{probbox}{41}{\stars{1}}
Hình tam giác ABC có diện tích $12$ cm$^2$. Kéo dài AB về phía B một đoạn $BD=AB$ và kéo dài AC về phía C một đoạn $CE=2\times AC$. Tính diện tích hình tam giác ADE.
\end{probbox}
\probsol{
$AD=2\times AB$, $AE=3\times AC$.\par Nối D với C.\par ADC và ABC chung chiều cao hạ từ C: $S_{ADC}=12\times2=24$ cm$^2$.\par ADE và ADC chung chiều cao hạ từ D: $S_{ADE}=24\times3=72$ cm$^2$.
\dapso{$72$ cm$^2$.}}
\gap
\begin{probbox}{42}{\stars{1}}
Hình tam giác ABC có diện tích $15$ cm$^2$. Kéo dài AB về phía B một đoạn $BD=2\times AB$ và kéo dài AC về phía C một đoạn $CE=2\times AC$. Tính diện tích hình tam giác ADE.
\end{probbox}
\probsol{
$AD=3\times AB$, $AE=3\times AC$.\par Nối D với C.\par ADC và ABC chung chiều cao hạ từ C: $S_{ADC}=15\times3=45$ cm$^2$.\par ADE và ADC chung chiều cao hạ từ D: $S_{ADE}=45\times3=135$ cm$^2$.
\dapso{$135$ cm$^2$.}}
\gap
\begin{probbox}{43}{\stars{2}}
Hình tam giác ABC có diện tích $80$ cm$^2$. Các điểm M, N, P lần lượt nằm trên AB, BC, CA sao cho $AM=\dfrac{1}{4}AB$, $BN=\dfrac{1}{4}BC$, $CP=\dfrac{1}{4}CA$. Tính diện tích hình tam giác MNP.
\end{probbox}
\probsol{
Trong hình tam giác AMP: $AM=\dfrac{1}{4}AB$, $AP=\dfrac{3}{4}AC$.\par Làm như trên (nối M với C): $S_{AMP}=80\times\dfrac{1}{4}\times\dfrac{3}{4}=15$ cm$^2$.\par Tương tự, $S_{BNM}=S_{CPN}=15$ cm$^2$.\par $S_{MNP}=80-3\times15=35$ cm$^2$.
\dapso{$35$ cm$^2$.}}
\gap
\begin{probbox}{44}{\stars{2}}
Hình tam giác ABC có diện tích $72$ cm$^2$. Các điểm M, N, P lần lượt nằm trên AB, BC, CA sao cho $AM=\dfrac{1}{3}AB$, $BN=\dfrac{1}{3}BC$, $CP=\dfrac{1}{3}CA$. Tính diện tích hình tam giác MNP.
\end{probbox}
\probsol{
Trong hình tam giác AMP: $AM=\dfrac{1}{3}AB$, $AP=\dfrac{2}{3}AC$.\par Làm như trên (nối M với C): $S_{AMP}=72\times\dfrac{1}{3}\times\dfrac{2}{3}=16$ cm$^2$.\par Tương tự, $S_{BNM}=S_{CPN}=16$ cm$^2$.\par $S_{MNP}=72-3\times16=24$ cm$^2$.
\dapso{$24$ cm$^2$.}}
\gap
\begin{probbox}{45}{\stars{1}}
Trong hình tam giác ABC, điểm M trên BC có $BM:MC=3:1$. Diện tích hình tam giác ABM lớn hơn diện tích hình tam giác AMC là $20$ cm$^2$. Tính diện tích hình tam giác ABC.
\end{probbox}
\probsol{
ABM và AMC chung chiều cao hạ từ A, nên $S_{ABM}:S_{AMC}=3:1$.\par Hiệu $2$ phần ứng với $20$ cm$^2$; một phần là $10$ cm$^2$.\par $S_{ABC}=10\times4=40$ cm$^2$.
\dapso{$40$ cm$^2$.}}
\gap
\begin{probbox}{46}{\stars{1}}
Trong hình tam giác ABC, điểm M trên BC có $BM:MC=5:2$. Diện tích hình tam giác ABM lớn hơn diện tích hình tam giác AMC là $24$ cm$^2$. Tính diện tích hình tam giác ABC.
\end{probbox}
\probsol{
ABM và AMC chung chiều cao hạ từ A, nên $S_{ABM}:S_{AMC}=5:2$.\par Hiệu $3$ phần ứng với $24$ cm$^2$; một phần là $8$ cm$^2$.\par $S_{ABC}=8\times7=56$ cm$^2$.
\dapso{$56$ cm$^2$.}}
\gap
\begin{probbox}{47}{\stars{1}}
Hình chữ nhật ABCD có $AB=24$ cm, $BC=16$ cm. Điểm M trên AB có $AM=8$ cm; điểm N trên BC có $BN=6$ cm. Tính diện tích hình tam giác DMN.
\end{probbox}
\probsol{
$S_{ABCD}=24\times16=384$ cm$^2$.\par $S_{AMD}=64$, $S_{MBN}=48$, $S_{NCD}=120$ cm$^2$.\par $S_{DMN}=384-64-48-120=152$ cm$^2$.
\dapso{$152$ cm$^2$.}}
\gap
\begin{probbox}{48}{\stars{1}}
Hình chữ nhật ABCD có $AB=30$ cm, $BC=20$ cm. Điểm M trên AB có $AM=12$ cm; điểm N trên BC có $BN=10$ cm. Tính diện tích hình tam giác DMN.
\end{probbox}
\probsol{
$S_{ABCD}=30\times20=600$ cm$^2$.\par $S_{AMD}=120$, $S_{MBN}=90$, $S_{NCD}=150$ cm$^2$.\par $S_{DMN}=600-120-90-150=240$ cm$^2$.
\dapso{$240$ cm$^2$.}}
\gap
\begin{probbox}{49}{\stars{1}}
Hình chữ nhật ABCD có $AB=18$ cm, $BC=12$ cm. Điểm M trên AB có $AM=6$ cm; điểm N trên BC có $BN=4$ cm. Tính diện tích hình tam giác DMN.
\end{probbox}
\probsol{
$S_{ABCD}=18\times12=216$ cm$^2$.\par $S_{AMD}=36$, $S_{MBN}=24$, $S_{NCD}=72$ cm$^2$.\par $S_{DMN}=216-36-24-72=84$ cm$^2$.
\dapso{$84$ cm$^2$.}}
\gap
\begin{probbox}{50}{\stars{1}}
Hình tam giác ABC có $BC=16$ cm. Kéo dài BC về phía C thêm $5$ cm thì được một hình tam giác có diện tích lớn hơn diện tích hình tam giác ABC là $25$ cm$^2$. Tính diện tích hình tam giác ABC.
\end{probbox}
\probsol{
Phần tăng thêm là hình tam giác có đáy $5$ cm và có cùng chiều cao hạ từ A.\par Chiều cao: $25\times2:5=10$ cm.\par $S_{ABC}=16\times10:2=80$ cm$^2$.
\dapso{$80$ cm$^2$.}}
\clearpage
\titlehead
